\section{Literature review}

This chapter provides an overview of the existing research and establishes the theoretical foundation for the work presented in this thesis.

\subsection{Background}

The field under investigation has seen significant developments over the past decade.
Early approaches focused on manual processes, which were time-consuming and error-prone \cite{sample-design}.
As the field matured, automated solutions emerged that addressed many of these limitations.

Table~\ref{tab:comparison} presents a comparison of existing approaches.

% === EXAMPLE: Basic table with horizontal and vertical lines ===
\begin{table}[H]
    \caption{Comparison of existing approaches} \label{tab:comparison}
    \begin{tabular}{|l|c|c|c|}
    \hline \textbf{Approach} & \textbf{Accuracy (\%)} & \textbf{Speed (ms)} & \textbf{Scalability} \\
    \hline Method A & 85 & 120 & Low \\
    \hline Method B & 90 & 95  & Medium \\
    \hline Method C & 88 & 110 & High \\
    \hline Method D & 92 & 80  & Medium \\
    \hline
    \end{tabular}
\end{table}

\subsection{Theoretical framework}

The theoretical framework for this research draws from several disciplines, including computer science, statistics, and domain-specific knowledge.

% === EXAMPLE: Numbered equation with sum ===
The following equation describes the fundamental relationship:

\begin{equation} \label{eq:basic}
    E = \sum_{i=1}^{n} w_i \cdot f(x_i)
\end{equation}

where $E$ is the overall effectiveness metric, $w_i$ are the weight coefficients, and $f(x_i)$ represents the performance function for each component.

% === EXAMPLE: Inline math ===
The time complexity of the algorithm is $O(n \log n)$ in the average case and $O(n^2)$ in the worst case.
The space complexity is bounded by $\Theta(n)$.

% === EXAMPLE: Multi-line aligned equations ===
The system of equations describing the optimization problem:

\begin{align}
    \min_{x} \quad & f(x) = \sum_{i=1}^{n} c_i x_i \label{eq:objective} \\
    \text{s.t.} \quad & \sum_{i=1}^{n} a_{ij} x_i \leq b_j, \quad j = 1, \ldots, m \label{eq:constraint1} \\
    & x_i \geq 0, \quad i = 1, \ldots, n \label{eq:constraint2}
\end{align}

% === EXAMPLE: Equation with fractions, integrals, and limits ===
The probability density function of the normal distribution is given by:

\begin{equation} \label{eq:normal}
    f(x \mid \mu, \sigma^2) = \frac{1}{\sigma\sqrt{2\pi}} \exp\left( -\frac{(x - \mu)^2}{2\sigma^2} \right)
\end{equation}

The cumulative distribution function is:

\begin{equation} \label{eq:cdf}
    F(x) = \int_{-\infty}^{x} f(t) \, dt = \frac{1}{2} \left[ 1 + \operatorname{erf}\left( \frac{x - \mu}{\sigma\sqrt{2}} \right) \right]
\end{equation}

% === EXAMPLE: Matrix notation ===
The transformation can be represented in matrix form:

\begin{equation} \label{eq:matrix}
    \mathbf{y} = \mathbf{A} \mathbf{x} + \mathbf{b}, \quad
    \text{where} \quad
    \mathbf{A} = \begin{pmatrix}
        a_{11} & a_{12} & \cdots & a_{1n} \\
        a_{21} & a_{22} & \cdots & a_{2n} \\
        \vdots & \vdots & \ddots & \vdots \\
        a_{m1} & a_{m2} & \cdots & a_{mn}
    \end{pmatrix}
\end{equation}

% === EXAMPLE: Cases / piecewise function ===
The activation function used in the neural network:

\begin{equation} \label{eq:relu}
    \operatorname{ReLU}(x) =
    \begin{cases}
        x & \text{if } x > 0 \\
        0 & \text{otherwise}
    \end{cases}
\end{equation}

\subsection{Related work}

Several researchers have contributed to this field.
Smith and Johnson \cite{sample-ref} proposed an approach based on modular architecture that achieved promising results.
Brown et al. \cite{sample-analysis} extended this work by introducing optimization techniques that improved performance by approximately 15\%.
Lee and Park \cite{sample-testing} conducted a comprehensive evaluation of multiple methods and identified key factors affecting system performance.

% === EXAMPLE: Nested enumerated list ===
The key contributions of prior work can be categorized as follows:
\begin{enumerate}
    \item \textbf{Architectural contributions:}
    \begin{enumerate}
        \item modular design patterns for distributed systems;
        \item event-driven architectures for real-time processing;
        \item containerized deployment strategies.
    \end{enumerate}
    \item \textbf{Algorithmic contributions:}
    \begin{enumerate}
        \item optimization heuristics for resource allocation;
        \item machine learning pipelines for automated classification;
        \item compression algorithms for efficient data transfer.
    \end{enumerate}
    \item \textbf{Evaluation methodologies:}
    \begin{enumerate}
        \item standardized benchmarking frameworks;
        \item statistical significance testing protocols;
        \item reproducibility guidelines.
    \end{enumerate}
\end{enumerate}

\subsection{Standards and best practices}

Industry standards play an important role in ensuring quality and interoperability.
Organizations such as \hyperlink{ieee}{IEEE}, \hyperlink{iso}{ISO}, and \hyperlink{nist}{NIST} have published guidelines that inform the design of modern systems \cite{nist-guide}.

% === EXAMPLE: TikZ diagram — service-oriented architecture ===
Figure~\ref{fig:architecture} illustrates a typical system architecture based on these standards.

\begin{figure}[H]
    \centering
    \begin{tikzpicture}[
        box/.style={rectangle, draw, minimum width=2.5cm, minimum height=0.8cm, align=center},
        >=Stealth
    ]
        \node[box] (ui) {User Interface};
        \node[box, below=0.6cm of ui] (api) {API Gateway};
        \node[box, below left=0.6cm and 0.3cm of api] (svc1) {Service A};
        \node[box, below right=0.6cm and 0.3cm of api] (svc2) {Service B};
        \node[box, below=0.6cm of api] (db) {Database};
        \draw[->] (ui) -- (api);
        \draw[->] (api) -- (svc1);
        \draw[->] (api) -- (svc2);
        \draw[->] (svc1) -- (db);
        \draw[->] (svc2) -- (db);
    \end{tikzpicture}
    \caption{General system architecture} \label{fig:architecture}
\end{figure}

% === EXAMPLE: TikZ flowchart with decision diamond ===
Figure~\ref{fig:flowchart} shows the general decision-making process described in the literature.

\begin{figure}[H]
    \centering
    \begin{tikzpicture}[
        startstop/.style={rectangle, rounded corners, draw, minimum width=2.5cm, minimum height=0.7cm, align=center, fill=blue!10},
        process/.style={rectangle, draw, minimum width=2.5cm, minimum height=0.7cm, align=center},
        decision/.style={diamond, draw, minimum width=2cm, minimum height=1cm, align=center, aspect=2},
        arrow/.style={->, >=Stealth, thick}
    ]
        \node[startstop] (start) {Start};
        \node[process, below=0.5cm of start] (input) {Load input data};
        \node[decision, below=0.7cm of input] (check) {Data valid?};
        \node[process, below left=0.7cm and 1.5cm of check] (error) {Report error};
        \node[process, below right=0.7cm and 1.5cm of check] (proc) {Process data};
        \node[process, below=0.7cm of proc] (output) {Generate report};
        \node[startstop, below=0.5cm of output] (stop) {End};

        \draw[arrow] (start) -- (input);
        \draw[arrow] (input) -- (check);
        \draw[arrow] (check) -- node[above] {No} (error);
        \draw[arrow] (check) -- node[above] {Yes} (proc);
        \draw[arrow] (proc) -- (output);
        \draw[arrow] (output) -- (stop);
        \draw[arrow] (error) |- (stop);
    \end{tikzpicture}
    \caption{Decision flowchart for data processing} \label{fig:flowchart}
\end{figure}

\subsection{Gap analysis}

Despite significant progress, several gaps remain in the current body of knowledge:
\begin{itemize}
  \item existing methods do not adequately address scalability;
  \item performance optimization under resource constraints is underexplored;
  \item integration with modern deployment pipelines requires further investigation.
\end{itemize}

This thesis addresses these gaps by proposing a novel approach that combines the strengths of existing methods while mitigating their weaknesses.

\clearpage
